\documentclass[aspectratio=169,11pt]{beamer}

% ============================================
% PACKAGES
% ============================================
\usepackage[utf8]{inputenc}
\usepackage[T1]{fontenc}
\usepackage[ngerman]{babel}
\usepackage{lmodern}
\usepackage{tcolorbox}
\tcbuselibrary{skins, hooks}
\usepackage{listings}
\usepackage{xcolor}
\usepackage{fontawesome5}
\usepackage{tikz}
\usetikzlibrary{shapes, arrows.meta, positioning}

% Modern Font (if installed, otherwise falls back smoothly)
\IfFileExists{FiraSans.sty}{\usepackage[sfdefault]{FiraSans}}{}
\renewcommand{\familydefault}{\sfdefault}

% ============================================
% COLORS & DESIGN SYSTEM
% ============================================
\definecolor{BgColor}{HTML}{FAFAFA}
\definecolor{Primary}{HTML}{2B3A42}
\definecolor{Accent}{HTML}{E84A5F}
\definecolor{Secondary}{HTML}{3F5765}
\definecolor{CodeBg}{HTML}{F0F0F0}
\definecolor{Success}{HTML}{28A745}

\setbeamercolor{background canvas}{bg=BgColor}
\setbeamercolor{title}{fg=Primary}
\setbeamercolor{frametitle}{bg=Primary, fg=white}
\setbeamercolor{structure}{fg=Accent}
\setbeamercolor{normal text}{fg=Primary}
\setbeamercolor{itemize item}{fg=Accent}
\setbeamercolor{itemize subitem}{fg=Secondary}

\setbeamertemplate{navigation symbols}{}
\setbeamertemplate{footline}{%
  \leavevmode%
  \hbox{%
  \begin{beamercolorbox}[wd=.333333\paperwidth,ht=2.25ex,dp=1ex,center]{author in head/foot}%
    \usebeamerfont{author in head/foot}\tiny Falko Himmel
  \end{beamercolorbox}%
  \begin{beamercolorbox}[wd=.333333\paperwidth,ht=2.25ex,dp=1ex,center]{title in head/foot}%
    \usebeamerfont{title in head/foot}\tiny Grundlagen der Praktischen Informatik
  \end{beamercolorbox}%
  \begin{beamercolorbox}[wd=.333333\paperwidth,ht=2.25ex,dp=1ex,right]{date in head/foot}%
    \usebeamerfont{date in head/foot}\insertframenumber{} / \inserttotalframenumber\hspace*{2ex} 
  \end{beamercolorbox}}%
  \vskip0pt%
}

% ============================================
% CUSTOM BOXES
% ============================================
\tcbset{
    modernbox/.style={
        enhanced,
        boxrule=0pt,
        frame hidden,
        colback=white,
        coltitle=white,
        colbacktitle=Secondary,
        fonttitle=\bfseries,
        drop lifted shadow,
        arc=4pt,
        toptitle=2mm,
        bottomtitle=2mm,
        left=2mm,
        right=2mm
    },
    alertbox/.style={
        modernbox,
        colbacktitle=Accent
    },
    successbox/.style={
        modernbox,
        colbacktitle=Success
    }
}

% ============================================
% CODE LISTINGS
% ============================================
\lstdefinelanguage{Haskell}{
  keywords={class, data, deriving, do, else, if, in, import, let, module, of, then, type, where, instance, func, return},
  keywordstyle=\color{Accent}\bfseries,
  ndkeywords={print, main},
  ndkeywordstyle=\color{Secondary}\bfseries,
  identifierstyle=\color{Primary},
  sensitive=true,
  comment=[l]{\#},
  commentstyle=\color{gray}\ttfamily\itshape,
  stringstyle=\color{teal}\ttfamily,
  morestring=[b]"
}

\lstset{
    language=Haskell,
    backgroundcolor=\color{CodeBg},
    basicstyle=\ttfamily\small,
    breaklines=true,
    frame=single,
    rulecolor=\color{CodeBg},
    frameround=tttt,
    numbers=left,
    numberstyle=\tiny\color{gray},
    xleftmargin=15pt,
    tabsize=4,
    showstringspaces=false
}

% ============================================
% TITLE PAGE SETTINGS
% ============================================
\title{\textbf{Tutorium 1}}
\subtitle{Endrekursion vs. Lineare Rekursion}
\institute{Grundlagen der Praktischen Informatik}
\author{}
\date{\today}

\begin{document}

\begin{frame}[plain]
    \titlepage
\end{frame}

\begin{frame}{Inhaltsverzeichnis}
    \tableofcontents
\end{frame}

% ============================================
% THEORIE
% ============================================
\section{Theorie: Rekursionsarten}

\begin{frame}{Klausuraufgabe}
    \begin{tcolorbox}[modernbox, title=\faQuestionCircle\ Typische Frage]
        Bewerten Sie, ob die folgende Aussage korrekt ist und begründen Sie Ihre Antwort:
        \vspace{0.3cm}
        
        \textit{,,Jede endrekursive Funktion ist auch linear rekursiv umsetzbar.''}
    \end{tcolorbox}
    
    \vspace{0.5cm}
    \pause
    \begin{itemize}
        \item[\faCheckCircle] \textbf{Kurzantwort:} Ja!
        \item Jede endrekursive Funktion lässt sich in eine (äquivalente) linear rekursive Funktion umschreiben.
    \end{itemize}
\end{frame}

\begin{frame}[fragile]{Endrekursion $\rightarrow$ Lineare Rekursion}
    \textbf{Endrekursion:} Der rekursive Aufruf ist der letzte Ausdruck. Kein weiterer Work nach der Rückkehr!
    
    \begin{columns}[T]
        \begin{column}{0.48\textwidth}
            \begin{tcolorbox}[modernbox, title=Endrekursiv (mit Akkumulator)]
\begin{lstlisting}
func fak(n, acc):
    if n == 0:
        return acc
    return fak(n-1, n * acc)
\end{lstlisting}
            \end{tcolorbox}
            \small \textbf{Aufruf:} \texttt{fak(3, 1) $\rightarrow$ fak(2, 3) $\rightarrow$ ... $\rightarrow$ 6}
        \end{column}
        \begin{column}{0.48\textwidth}
            \begin{tcolorbox}[successbox, title=Linear Rekursiv]
\begin{lstlisting}
func fak(n):
    if n == 0:
        return 1
    return n * fak(n-1)
\end{lstlisting}
            \end{tcolorbox}
            \small \textbf{Aufruf:} \texttt{3 * fak(2) $\rightarrow$ 3 * 2 * fak(1) $\rightarrow$ ... $\rightarrow$ 6}
        \end{column}
    \end{columns}
\end{frame}

\begin{frame}[fragile]{Visueller Vergleich: Call Stack}
    \begin{columns}[T]
        \begin{column}{0.48\textwidth}
            \textbf{Endrekursiv (Tail Call)} \\
            \small Der Zustand wird im \texttt{acc} weitergegeben. Kein wachsender Stack!
            \vspace{0.3cm}
            \begin{center}
            \begin{tikzpicture}[node distance=0.6cm, every node/.style={draw, rectangle, minimum width=2.5cm, fill=BgColor, rounded corners=2pt}]
                \node (call1) {\texttt{fak(3, 1)}};
                \node[below=of call1] (call2) {\texttt{fak(2, 3)}};
                \node[below=of call2] (call3) {\texttt{fak(1, 6)}};
                \node[below=of call3] (call4) {\texttt{fak(0, 6)}};
                \node[draw=none, fill=none, below=of call4, font=\bfseries\color{Success}] (res) {\texttt{Return 6}};
                
                \draw[-{Latex}, Accent, thick] (call1) -- (call2);
                \draw[-{Latex}, Accent, thick] (call2) -- (call3);
                \draw[-{Latex}, Accent, thick] (call3) -- (call4);
                \draw[-{Latex}, Success, thick] (call4) -- (res);
            \end{tikzpicture}
            \end{center}
        \end{column}
        \begin{column}{0.48\textwidth}
            \textbf{Linear Rekursiv} \\
            \small Jeder Aufruf wartet auf das Ergebnis. Der Stack wächst!
            \vspace{0.3cm}
            \begin{center}
            \begin{tikzpicture}[node distance=0.6cm, every node/.style={draw, rectangle, minimum width=2.5cm, fill=BgColor, rounded corners=2pt}]
                \node (call1) {\texttt{3 * fak(2)}};
                \node[below=of call1] (call2) {\texttt{2 * fak(1)}};
                \node[below=of call2] (call3) {\texttt{1 * fak(0)}};
                \node[below=of call3] (call4) {\texttt{Return 1}};
                
                \draw[-{Latex}, Accent, thick] (call1.south) -- (call2.north) node[midway, right, draw=none, fill=none, font=\tiny, color=gray] {wartet};
                \draw[-{Latex}, Accent, thick] (call2.south) -- (call3.north);
                \draw[-{Latex}, Accent, thick] (call3.south) -- (call4.north);
                
                \draw[dashed, -{Latex}, Success, thick] (call4.east) to[bend right=45] node[right, draw=none, fill=none, font=\tiny] {1 * 1} (call3.east);
                \draw[dashed, -{Latex}, Success, thick] (call3.east) to[bend right=45] node[right, draw=none, fill=none, font=\tiny] {2 * 1} (call2.east);
                \draw[dashed, -{Latex}, Success, thick] (call2.east) to[bend right=45] node[right, draw=none, fill=none, font=\tiny] {3 * 2} (call1.east);
            \end{tikzpicture}
            \end{center}
        \end{column}
    \end{columns}
\end{frame}

\begin{frame}[fragile]{Lineare Rekursion $\rightarrow$ Endrekursion}
    Die Umwandlung von linear zu endrekursiv ist \textbf{nicht immer trivial}. Es hängt davon ab, ob das Ergebnis schrittweise als Akkumulator aufgebaut werden kann.
    
    \begin{tcolorbox}[alertbox, title=Beispiel: $f(x) = 2 \cdot f(x-1) + 1$]
\begin{lstlisting}
func f(x):
    if x == 0:
        return 1
    return 2 * f(x-1) + 1
\end{lstlisting}
    \end{tcolorbox}
    \vspace{0.2cm}
    \small \textbf{Problem:} Das Ergebnis ist auf jeder Ebene eine lineare Kombination. Eine direkte Endrekursion ohne Hilfsmittel (geschlossene Form oder iterative Akkumulatoren) ist nicht offensichtlich.
\end{frame}

\begin{frame}[fragile]{Lösung über die geschlossene Form}
    Oft hilft eine Umformung (geschlossene Form). Für unser Beispiel gilt:
    \[ f(x) = 2^{x+1} - 1 \]
    
    \begin{tcolorbox}[modernbox, title=Iterative/Endrekursive Umsetzung]
\begin{lstlisting}
func f_tail(x, acc):
    if x == 0:
        return acc
    return f_tail(x-1, 2*acc + 1)
\end{lstlisting}
    \end{tcolorbox}
    \small
    \textbf{Aufruf für $x=3, acc=1$:} \\
    \texttt{f\_tail(3, 1) $\rightarrow$ f\_tail(2, 3) $\rightarrow$ f\_tail(1, 7) $\rightarrow$ f\_tail(0, 15) $\rightarrow$ 15}
\end{frame}

\begin{frame}{Zusammenfassung}
    \begin{tcolorbox}[modernbox, title=\faListUl\ Wichtige Erkenntnisse]
        \begin{itemize}
            \item \textbf{Endrekursion $\rightarrow$ Linear:} Immer möglich! (Akkumulator entfernen, Rechnung beim Zurückkehren durchführen).
            \item \textbf{Linear $\rightarrow$ Endrekursion:} Nicht immer trivial. Manche linearen Rekursionen benötigen zusätzliche Akkumulatoren oder eine geschlossene Form.
            \item \textbf{Prinzip:} Endrekursion trägt das Zwischenergebnis in Parametern; lineare Rekursion baut das Ergebnis beim Zurückkehren auf.
        \end{itemize}
    \end{tcolorbox}
\end{frame}

% ============================================
% LIVE DEMO
% ============================================
\section{Praxis}

\begin{frame}[plain]
    \begin{center}
        \Huge \textbf{Live Demo} \\
        \vspace{0.5cm}
        \large \faLaptopCode~ \texttt{ex00.hs} \\
        \vspace{0.3cm}
        \normalsize Wir programmieren den \textbf{Hello World} und \textbf{ggT} (Größter gemeinsamer Teiler) zusammen im Code!
    \end{center}
\end{frame}

\end{document}
